\subsection{An exact parity formula for supersingular elliptic products}
\label{sec:DE-A3-ss-product}

The following family gives an explicit answer in every degree, and shows that
eventual relative ordinarity can fail even when the ambient Picard scheme is
smooth and its crystalline cohomology is torsion-free.  It does not purport to
determine the generic rank for an arbitrary fixed ambient variety.

Throughout this subsection let $k=\overline{\mathbf F}_2$, let
\[
 E_0:\quad y^2+y=x^3,
 \qquad X=E_0\times E_0,
 \qquad L=\OO_{E_0}(3O)\boxtimes\OO_{E_0}(3O),
\]
and embed $X\hookrightarrow\mathbf P^8_k$ by the product of the plane cubic
embeddings followed by the Segre embedding.  Write
$N_m=\OO_{E_0}(mO)\boxtimes\OO_{E_0}(mO)$, so $L^d=N_{3d}$.
For a smooth member $Y\in|N_m|$, put
\[
 Q_Y=\coker\bigl(H^1(X,\OO_X)\longrightarrow H^1(Y,\OO_Y)\bigr).
\]
Frobenius on these vector spaces is $a\mapsto a^2$-semilinear; ``stable rank''
means the dimension of the eventual image of its iterates over the perfect
field $k$.

\begin{lemma}[Explicit endomorphisms of every positive degree]
\label{lem:DE-A3-ss-norm}
\status{Proved}
For every integer $n\ge1$ there is an explicitly specified endomorphism
$b:E_0\to E_0$ with
$b^\dagger b=bb^\dagger=[n]$, where $\dagger$ is the Rosati involution for
$\OO_{E_0}(O)$.  In particular $\deg b=n$.
\end{lemma}
\begin{proof}
Choose $\omega\in\mathbf F_4$ with $\omega^2+\omega+1=0$, and define
\[
 i(x,y)=(x+1,y+x+\omega),\qquad
 j(x,y)=(x+\omega,y+\omega^2x+\omega).
\]
For $a\in\mathbf F_4^\times$ and $c^2+c=1$, substitution shows that
$(x,y)\mapsto(x+a,y+a^2x+c)$ preserves the equation, because $a^4=a$ and
$a^3=1$.  These maps extend to automorphisms fixing $O$.  Direct composition
gives $i^2=j^2=[-1]$ and $ij=-ji$, with $[-1](x,y)=(x,y+1)$.
Since they fix the polarizing divisor $O$, their Rosati adjoints are their
inverses.  Consequently
\[
 (a+bi+cj+dij)^\dagger(a+bi+cj+dij)
 =[a^2+b^2+c^2+d^2]
\]
for integers $a,b,c,d$; the same identity holds in the opposite order.

We recall a short proof that every positive integer is a sum of four squares,
so the construction does not depend on an unproved norm-surjectivity claim.
The quaternion norm identity makes products of such sums sums of four
squares.  For an odd prime $\ell$, the two sets of residues $a^2$ and
$-1-b^2$ each have $(\ell+1)/2$ elements.  They intersect.  Choosing
$|a|,|b|\le(\ell-1)/2$ gives
$a^2+b^2+1=k\ell$ with $0<k<\ell$.
Choose the smallest positive multiplier $k$ for which $k\ell$ is a sum of
four squares, and suppose $k>1$.  Write $k\ell=\sum a_i^2$ and choose
$x_i\equiv a_i\pmod k$ with $|x_i|\le k/2$.  Then
$\sum x_i^2=kt$ with $0\le t\le k$.  The case $t=0$ would imply
$k\mid\ell$, which is impossible.  If $t<k$, all four coefficients of the
quaternion product $a\overline x$ are divisible by $k$, since modulo $k$
this product is $x\overline x=kt$.  Dividing those coefficients by $k$
expresses $t\ell$ as a sum of four squares, contradicting minimality.
If $t=k$, equality in the four bounds forces $k$ even and every
$a_i\equiv k/2\pmod k$.  Hence $k\mid4\ell$, so $k=2$ or $4$.
For $k=2$, the four odd squares have sum $4\pmod8$, whereas $2\ell$ is
$2$ or $6\pmod8$.  For $k=4$, the four squares are each $4\pmod{16}$,
whereas $4\ell$ is $4$ or $12\pmod{16}$.  Both are impossible.  Thus
$k=1$.  The prime $2$ and the integer $1$ are immediate, and multiplicativity
finishes the proof.  To specify $b$ without choices, take the first
nonnegative quadruple in lexicographic order with sum of squares $n$ and
use the displayed endomorphism.  Finally
$(\deg b)^2=\deg[n]=n^2$, so $\deg b=n$.
\end{proof}

\begin{lemma}[The ambient quotient and the even-degree obstruction]
\label{lem:DE-A3-ss-prym}
\status{Proved}
For every smooth $Y\in|N_m|$, $m\ge3$, its genus is $g=m^2+1$ and
$\dim Q_Y=m^2-1$.  Its $2$-rank $f_Y$ satisfies
\[
 f_Y\le m^2-1,
 \qquad
 f_Y\le m^2-3\quad\hbox{if }2\mid m.
\]
Moreover $\dim(Q_Y)_{\ss}=f_Y$.
\end{lemma}
\begin{proof}
The plane cubic is smooth: on the affine chart the partial derivative with
respect to $y$ is $1$, and at its unique point at infinity the homogeneous
equation has a nonzero partial derivative.  Its inverse map is
$(x,y)\mapsto(x,y+1)$, so $E_0[2](k)=0$.  Thus its $2$-rank is zero and
Frobenius on $H^1(E_0,\OO_{E_0})$ is zero.  The latter also follows from
Cartier: $\omega_0=dx$ is a regular invariant differential and
$C(dx)=0$.  Hence Frobenius is zero on $H^1(X,\OO_X)$, a space of
dimension two.

K\"unneth gives $H^1(X,P\otimes N_m^{-1})=0$ for every
$P\in\Pic^0(X)$: on each elliptic factor the two line bundles in question
have negative degree, so both possible $H^1$ tensor summands have a zero
$H^0$ factor.  In particular $H^0(Y,\OO_Y)=k$ and
$H^1(X,\OO_X)\hookrightarrow H^1(Y,\OO_Y)$.  Adjunction, or the same
exact sequence and Serre duality, gives $g=m^2+1$.

Restriction
\[
 j:\Pic^0(X)=E_0^2\longrightarrow J_Y
\]
is a closed immersion of group schemes.  Indeed a nontrivial
$P\in\Pic^0(X)(k)$ has $H^0(X,P)=0$.  The preceding vanishing and the
restriction sequence then give $H^0(Y,P|_Y)=0$, so $P|_Y$ is not trivial.
The tangent map of $j$ is the already proved injection on $H^1(\OO)$.
Its proper kernel therefore has only the identity as a geometric point and
has zero tangent space.  It is a finite local group scheme with zero
cotangent space, and Nakayama's lemma makes it trivial.

Identify $X$ with its dual by the product principal polarization.  The map
dual to $j$ is
\[
 q:J_Y\longrightarrow X,\qquad
 M\longmapsto
 \bigl(\operatorname{Nm}_{\pi_1}M,
       \operatorname{Nm}_{\pi_2}M\bigr),
\]
where $\pi_i:Y\to E_0$ are the projections.  They are finite of degree
$m$: $Y$ is smooth and connected, hence integral, and cannot contain an
elliptic fibre.  The diagonal entries of $qj$ are $[m]$ by norm and
pullback.  Each off-diagonal entry is zero.  For example the intersection
divisor on the first $E_0$ obtained by fixing a point on the second factor
always belongs to $|\OO_{E_0}(mO)|$, so the norms of the pullbacks of two
point divisors on the second factor are linearly equivalent.  We obtain
\[
 qj=[m]_X.                                                     \tag{\ast}
\]
In particular $q$ is surjective.

For completeness its kernel is connected and smooth, not merely connected
up to isogeny.  Let $P=(\ker q)^0_{\mathrm{red}}$.  Over the perfect field
$k$ this is an abelian subvariety.  The quotient $J_Y/P$ is an abelian
variety and $q$ factors through an isogeny $\nu:J_Y/P\to X$.
Dualizing, the kernel of $\nu^t$ lies in the kernel of $q^t=j$, which is
trivial.  The isogeny duality theorem gives
$\ker(\nu^t)\simeq(\ker\nu)^D$, so $\nu$ is an isomorphism.
Thus $\ker q=P$.  Here the exact source statements are the quotient
existence theorem and its abelian-variety example, and the isogeny duality
theorem \cite[Theorem~4.38, Example~4.40, Theorem~7.5]{EdixhovenVanDerGeerMoonenAV}.

The $2$-rank is additive in an exact sequence of abelian varieties.
One can see this without choosing an isogeny complement: on $k$-points,
the sequence of $2^a$-torsion groups is exact, since a lift of a torsion
point can be adjusted by a point of $P(k)$ using the surjectivity of
$[2^a]:P(k)\to P(k)$.  Counting points gives additivity.
Since $f_X=0$, we have $f_Y=f_P$ and $\dim P=g-2$.  This proves the first
bound.

If $m$ is even, (\ast) restricts to a closed immersion of finite group
schemes $j_2:X[2]\hookrightarrow P[2]$.  Its Frobenius twist maps
$X^{(2)}[2]$ into $P^{(2)}[2]$.  The group $\ker V_X$ lies in
$X^{(2)}[2]$, and finite-group Verschiebung is functorial for $j_2$
(or use $V_{J_Y}j^{(2)}=jV_X$ and the restriction of $V_{J_Y}$ to
$P^{(2)}$).  Consequently it embeds
\[
 \ker(V_X:X^{(2)}\to X)\lhook\joinrel\longrightarrow
 \ker(V_P:P^{(2)}\to P).
\]
The source is connected of order $4$, since $\deg V_X=2^2$ and $f_X=0$.
For any abelian variety $A$ of dimension $t$, $\deg V_A=2^t$ and
$\#\ker V_A(k)=2^{f_A}$: relative Frobenius is bijective on $k$-points
and $V_AF_A=[2]$.  The identity component of $\ker V_A$ therefore has
order $2^{t-f_A}$.  The displayed injection yields
$\dim P-f_P\ge2$, hence $f_Y\le g-4=m^2-3$.
These degree and Frobenius identities are exactly
\cite[Propositions~5.19--5.22]{EdixhovenVanDerGeerMoonenAV}; the argument
uses the connected subgroup, so it is not just a bound on geometric
$2$-torsion points.

Finally Artin--Schreier and the semilinear Fitting decomposition identify
the stable rank of $H^1(Y,\OO_Y)$ with $f_Y$.  The stable-part functor on
finite-dimensional semilinear vector spaces over a perfect field is exact:
split each space into its nilpotent and bijective parts, or apply a
sufficiently high Frobenius iterate to an exact sequence.  Quotienting by
the Frobenius-zero space $H^1(X,\OO_X)$ therefore leaves the stable rank
unchanged.  This proves the last assertion.
\end{proof}

\begin{theorem}[Exact generic rank in every hypersurface degree]
\label{thm:DE-A3-ss-product-rank}
\status{Proved}
For the specified $X\subset\mathbf P^8_{\mathbf F_2}$, the geometric
generic stable rank in degree $d\ge1$ is
\[
 \boxed{\quad
 \dim E(Y,\OO_Y)_{\ss,\mathrm{gen}}
 =\begin{cases}
 9d^2-1,&d\ \text{odd},\\
 9d^2-3,&d\ \text{even}.
 \end{cases}\quad}                                           \tag{\dagger}
\]
The coherent relative space has dimension $9d^2-1$ in every degree.
Consequently its generic Frobenius is bijective exactly for odd $d$.
There is no integer $d_0$ making it generically bijective for every
$d\ge d_0$.  The generic stable rank is positive in every degree, so
for this example the positivity threshold is exactly $d_0=1$.
\end{theorem}
\begin{proof}
First let $m\ge3$ be odd.  Choose the endomorphism $b$ in
Lemma~\ref{lem:DE-A3-ss-norm} with $b^\dagger b=bb^\dagger=[m-1]$ and put
\[
 \phi=\begin{pmatrix}1&b\\-b^\dagger&1\end{pmatrix}:X\longrightarrow X.
\]
Its Rosati adjoint satisfies $\phi^\dagger\phi=[m]I$; hence $\phi$ is an
isogeny of degree $m^2$ with kernel killed by the odd integer $m$.
It is therefore \'etale.  The inverse images of the two coordinate
elliptic divisors are
\[
 C_1=\{x+by=0\},\qquad C_2=\{-b^\dagger x+y=0\}.
\]
Each is a graph isomorphic to $E_0$.  They meet transversely in
$\ker\phi$, which consists of exactly $m^2$ reduced points.

We need linear equivalence with the fixed $N_m$.  The polarization of
$\OO(C_1+C_2)=\phi^*(\OO_{E_0}(O)\boxtimes\OO_{E_0}(O))$ is $[m]I$,
by the pullback formula for polarizations
\cite[Proposition~7.6]{EdixhovenVanDerGeerMoonenAV}.
This is also the polarization of $N_m$, so their difference belongs to
$\Pic^0(X)$ by \cite[Corollary~7.22]{EdixhovenVanDerGeerMoonenAV}.  Both line bundles are symmetric; on $\Pic^0(X)$ inversion
pulls a line bundle back to its inverse.  Their difference is thus a
$2$-torsion $k$-point of $\Pic^0(X)$.  There are no such nontrivial points,
since $E_0[2](k)=0$.  Hence $C_1+C_2\in|N_m|$ exactly.

For a connected nodal curve $D$ with smooth components $D_i$, the
normalization sequence gives a Frobenius-equivariant exact sequence
\[
 0\longrightarrow H^1(\Gamma_D,k)
 \longrightarrow H^1(D,\OO_D)
 \longrightarrow\bigoplus_iH^1(D_i,\OO_{D_i})\longrightarrow0,
\]
where $\Gamma_D$ is its dual graph and Frobenius on the first term is
bijective.  Thus its stable rank is
$b_1(\Gamma_D)+\sum_i f_{D_i}$, and its genus minus stable rank is
$\sum_i(g(D_i)-f_{D_i})$.  Applied to $C_1\cup C_2$, this gives genus
$m^2+1$ and stable rank $m^2-1$.

The universal Cartier divisor over the projective space $|N_m|$ is proper
and flat.  Every nonzero section has connected divisor and
$h^1(\OO_D)=m^2+1$: this follows from
$H^1(N_m^{-1})=0$ and the surjection
$H^2(N_m^{-1})\to H^2(\OO_X)$ dual to the nonzero section.
Cohomology and base change therefore make $H^1(\OO_D)$ a vector bundle
on the parameter space.  Its relative Frobenius is a bundle map from the
Frobenius pullback.  For vector spaces of dimension $g=m^2+1$, the stable
rank is the rank of the $g$th iterate.  Rank at least a fixed integer is
an open condition.  The nodal member just constructed consequently
supplies a nonempty open locus of rank at least $m^2-1$.
Since $N_m$ is very ample for $m\ge3$, the smooth-member locus is also
nonempty open by Bertini.  These opens meet in the irreducible projective
space $|N_m|$.  Lemma~\ref{lem:DE-A3-ss-prym} proves that the generic
smooth member has rank exactly $m^2-1$.

Now let $m=3d$ be even, so $m\ge6$.  Write $m=3+(m-3)$, with both
summands odd and at least three.  Choose smooth members
$C\in|N_3|$ and $D\in|N_{m-3}|$ in the just proved maximal-rank opens.
They may be chosen transverse: both systems are very ample, so a general
member of the second system meets a fixed smooth member of the first
transversely.  Their union is a connected nodal member of $|N_m|$.
Each component has genus minus stable rank two.  The normalization
formula therefore gives
\[
 g(C\cup D)-f_{C\cup D}=4,\qquad f_{C\cup D}=m^2-3.
\]
The same open-rank and smoothness argument gives generic smooth rank at
least $m^2-3$.  The even-degree upper bound of
Lemma~\ref{lem:DE-A3-ss-prym} makes this an equality.

Finally these are precisely the complete linear systems cut out by
ambient hypersurfaces of degree $d$.  For the plane cubic,
$H^0(\mathbf P^2,\OO(d))\to H^0(E_0,\OO(3dO))$ is onto, by
$H^1(\mathbf P^2,\OO(d-3))=0$.  The Segre coordinate ring has all its
bihomogeneous pieces of bidegree $(d,d)$ generated by degree-$d$
polynomials in the pairwise products of coordinates.  Consequently
\[
 H^0(\mathbf P^8,\OO(d))\twoheadrightarrow H^0(X,L^d).
\]
The nonempty open rank and smoothness loci pull back to nonempty open
loci in the ambient hypersurface parameter space.  Taking $m=3d$ gives
(\dagger).  The entire construction is defined over an algebraic
extension of $\mathbf F_2$ and computes the geometric generic ranks for
the displayed variety over $\mathbf F_2$ itself.
\end{proof}

\begin{corollary}[An ambient Cartier-nilpotent dimension is insufficient]
\label{cor:DE-A3-ss-product-defect}
\status{Proved}
For this abelian surface, $\Pic^0(X)$ is smooth and
$H^*_{\cris}(X/W(k))$ is torsion-free.  All its global one-forms are
Cartier-killed, so the space of global forms eventually killed by Cartier
has dimension two.  Nevertheless the generic curve has
$g(Y)-f_Y=2$ in odd degree and $g(Y)-f_Y=4$ in even degree.
Thus neither eventual relative ordinarity nor a formula equating the
generic full defect with this ambient Cartier-nilpotent dimension holds
for every sufficiently large degree.
\end{corollary}
\begin{proof}
The Picard variety is the smooth abelian surface $E_0^2$.  Crystalline
cohomology of a product of elliptic curves is the exterior algebra on
their free first crystalline cohomology, hence is torsion-free
\cite[II, (7.1.1), p.~651]{Illusie1979DRW}.
The two pullbacks of $dx$ span the global one-forms and have Cartier zero.
The defect assertion is Theorem~\ref{thm:DE-A3-ss-product-rank} and
$g(Y)=9d^2+1$.
\end{proof}

\begin{remark}[Computations and exact scope]
\label{rem:DE-A3-ss-product-computation}
The exact finite-field probe
\texttt{A3\_supersingular\_product\_probe.py} uses the basis
$\{x^a:2a\le m\}\cup\{x^ay:2a+3\le m\}$ of
$H^0(E_0,\OO(mO))$ and the Cartier identities
\[
 C(x^{2a}dx)=0,\quad C(x^{2a+1}dx)=x^a dx,\quad
 C(x^{2a}y\,dx)=x^{a+1}dx,\quad
 C(x^{2a+1}y\,dx)=x^ay\,dx.
\]
Its recorded best stable matrix ranks for $d=1,\ldots,6$ are
$8,33,80,141,224,321$; elapsed time is
$29.180787964000046$ seconds.  It retains the exact section coefficients,
matrix columns, rank sequences and seeds.  These finite probes motivated
the parity formula; the nodal deformation and Prym arguments prove it
uniformly, without inferring either smoothness or an all-degree theorem
from random samples.  The additional certificate
\texttt{A3\_supersingular\_product\_checks.py} checks the explicit
automorphisms, both symbolic quaternion norm identities, four-square
witnesses for $1\le n\le4096$, the degeneration arithmetic for
$1\le d\le10000$, and all six retained matrix witnesses.  It passed
in $0.027462210000521736$ seconds; its complete code and output are retained.
The generic-rank problem for arbitrary $X$ in A.3 remains a separate
question.  No priority claim for this explicit formula is made here.
\end{remark}
